\section{模型假设与符号说明}

\subsection{模型假设}

\begin{enumerate}
  \item 题目提供的销售明细和外部变量记录总体真实可靠。对于已发现的数据问题，本文采用记录、标记和统一口径处理，而不是直接删除原始数据。
  \item 短期内历史销售规律具有一定延续性。未来 7 天预测依赖最近历史销量、同星期规律和类别/门店稳定差异。
  \item 负销量不代表顾客正常购买需求。根据人工确认，负销量多对应破损、过期或冲销类负向调整，因此预测目标采用正向销量。
  \item 附件类别字段可以作为同类零食的正式聚合依据。相关性聚类和销量规模分组只作为补充分析，不替代业务类别。
  \item 天气、节假日和活动日与销量之间可以存在统计关联，但当前数据不足以识别严格因果效应。本文所称“影响因素分析”均指统计关联分析。
  \item 未来 7 天无附件真实天气和真实活动安排。最终预测若使用天气和活动变量，应解释为“给定历史同期外部变量情景下”的条件预测。
\end{enumerate}

\subsection{符号说明}

\begin{table}[H]
\centering
\caption{主要符号说明}
\label{tab:symbols}
\begin{tabular}{ll}
\toprule
符号 & 含义 \\
\midrule
$t$ & 日期序号 \\
$s$ & 门店编号 \\
$p$ & 商品编号 \\
$g$ & 商品类别 \\
$y_{s,p,t}$ & 门店 $s$、商品 $p$ 在日期 $t$ 的正向销量 \\
$\hat y_{s,p,t}$ & 对 $y_{s,p,t}$ 的预测值 \\
$Y_{t,g}$ & 日期 $t$ 类别 $g$ 的聚合销量 \\
$S_t$ & 指数平滑中的平滑水平 \\
$\alpha$ & 指数平滑系数 \\
$r_{pq}$ & 商品 $p$ 与商品 $q$ 的销量相关系数 \\
$X_t$ & 日历、天气、活动日等外部变量 \\
$MAE$ & 平均绝对误差 \\
$RMSE$ & 均方根误差 \\
$WAPE$ & 加权绝对百分比误差 \\
\bottomrule
\end{tabular}
\end{table}
